\writeSectionFrame{\presentationTitle}{Fazit}
\begin{frame}{Anwendungsfälle}
	\begin{itemize}
 		\item Operationen haben verschiedene Zuständigkeiten
 		\item Operationen sind leicht zu erweitern
 		\item Schnittstellen der Elementklassen sollen schlank bleiben
 		\item Objektstruktur besitzt viele verschiedene Elemente
 		\item Man möchte nicht, dass die Elemente die Operationen selbst enthalten
 		\item Elementklassen bleiben konstant
 	\end{itemize}
\end{frame}

\realworld{Java Compiler API (javax.lang.model.element)}{Analyse des Quellcodes}
\note{
    In der Java Compiler API, die seit Java 6 verfügbar ist, wird das Besucher-Muster intensiv verwendet, um den Quellcode während der Kompilierung zu analysieren. Die API enthält das javax.lang.model.element.ElementVisitor-Interface, das verwendet wird, um verschiedene Elementtypen wie Klassen, Methoden und Felder zu besuchen und Operationen auf ihnen auszuführen.
}

\realworld{ANTLR (ANother Tool for Language Recognition)}{Traversieren von Abstrakten Syntaxbäumen}
\note{
    ANTLR ist ein weit verbreitetes Framework zum Generieren von Parsern und Lexer. Es verwendet das Besucher-Muster, um Abstrakte Syntaxbäume (ASTs) zu traversieren. Wenn ein AST erstellt wird, kann ein Visitor durch diesen Baum gehen, um spezifische Operationen wie Codegenerierung, Analyse oder Transformationen durchzuführen.
}

\realworld{Apache Tomcat}{Analyse der Struktur von JSP-Seiten}
\note{
    Apache Tomcat verwendet das Besucher-Muster, um die Struktur von JSP-Seiten zu analysieren und zu verarbeiten. JSP-Tags werden in eine interne Baumstruktur umgewandelt, und ein Visitor wird verwendet, um diese Struktur zu durchlaufen und die JSP-Seiten in Servlets zu übersetzen.
}

\realworld{XML Processing (JDOM)}{Durchlaufen der Baumstruktur eines XML-Dokuments}
\note{
    Im Bereich der XML-Verarbeitung wird das Besucher-Muster ebenfalls häufig verwendet. In Bibliotheken wie JDOM oder DOM4J wird es verwendet, um die Baumstruktur eines XML-Dokuments zu durchlaufen und verschiedene Operationen auszuführen, wie das Extrahieren von Daten, das Validieren des Dokuments oder das Konvertieren in andere Formate.
}

\realworld{Hibernate}{Objektmodell und Metadaten durchlaufen}
\note{
    In Hibernate wird das Besucher-Muster verwendet, um das Objektmodell und die dazugehörigen Metadaten zu durchlaufen, insbesondere im Zusammenhang mit dem Erstellen von SQL-Abfragen oder dem Durchführen von Validierungen.
}

\vorteil{Trennung von Struktur und Verhalten}
\note{
    Das Muster ermöglicht es, die Struktur von Objekten von den darauf angewendeten Operationen zu trennen. Dadurch können neue Operationen hinzugefügt werden, ohne die Klassenstruktur zu ändern
}

\vorteil{Einfache Erweiterbarkeit von Operationen}
\note{
    Neue Operationen können durch das Hinzufügen neuer Visitor-Klassen implementiert werden, ohne bestehende Element-Klassen zu modifizieren. Dies ist besonders nützlich, wenn viele verschiedene Operationen auf denselben Objekten ausgeführt werden müssen. Das Open-Closed-Prinzip ist also hinsichtlich der Operationen erfüllt.
}

\vorteil{Zentrale Verwaltung von Operationen}
\note{
    Das Verhalten wird in den Visitor-Klassen zentralisiert, was die Verwaltung und Pflege der Operationen erleichtert. Dies führt oft zu einer besseren Übersichtlichkeit des Codes.
}

\vorteil{Unterstützung von nicht verwandten Operationen}
\note{
    Das Muster ist nützlich, wenn unterschiedliche Operationen, die nicht eng miteinander verbunden sind, auf dieselben Objektstrukturen angewendet werden sollen. Jede Operation kann in einem eigenen Visitor implementiert werden.
}

\vorteil{Besucher kann wichtige Informationen sammeln}
\note{
    Mit dem Besucher gibt es auch eine einfache Möglichkeit, Informationen über alle Elemente zu sammeln und auszuwerten.
}

\nachteil{Schwierigkeit bei der Erweiterung der Objektstruktur}
\note{
    Das Hinzufügen neuer Element-Klassen (d.h. neuer Objekttypen) erfordert Änderungen an allen bestehenden Visitor-Klassen, was aufwendig und fehleranfällig sein kann. Das Muster ist daher weniger geeignet, wenn die Struktur der Objekte häufig geändert wird. Hinsichtlich der Elementtypen ist das Open-Closed-Prinzip also nicht erfüllt. 
}

\nachteil{Umständlichkeit bei der Implementierung}
\note{
    Das Implementieren eines Besuchers kann kompliziert werden, insbesondere wenn viele verschiedene Objekttypen und Besucher beteiligt sind. Dies führt oft zu einer erhöhten Komplexität und längeren Entwicklungszeiten.
}

\nachteil{Unabhängigkeit von der Objektstruktur}
\note{
    Das Besucher-Muster setzt voraus, dass alle beteiligten Objekte eine einheitliche Schnittstelle (Element) implementieren, was in manchen Fällen zu einer starren und weniger flexiblen Architektur führen kann.
}

\nachteil{Mögliche Verletzung des Kapselungsprinzips}
\note{
    Da das Besucher-Muster auf die internen Details der Objekte zugreift, besteht die Gefahr, dass es das Kapselungsprinzip verletzt, wenn Besucher zu viel Wissen über die internen Strukturen der Element-Klassen erfordern.
}

\nachteil{Verständlichkeit und Wartbarkeit}
\note{
    In großen Systemen mit vielen verschiedenen Besuchern und Elementen kann es schwierig werden, den Überblick über die verschiedenen Interaktionen zu behalten, was die Wartbarkeit des Systems beeinträchtigen kann.
}